% ==============================================================================
% CONFIGURACIÓN DE PAQUETES Y ESTILOS - LATEX
% UNSAAC | Asignatura: Ingeniería de Software I
% Laboratorio N° 05: Estimación de Esfuerzo con Planning Poker
% Proyecto: Sistema Integral de Gestión de Capacitaciones (SIGC-CUSCO)
% ==============================================================================

\usepackage[utf8]{inputenc}
\usepackage[T1]{fontenc}
\usepackage[spanish, es-tabla, es-nodecimaldot]{babel}
\spanishdeactivate{"}
\usepackage{csquotes}

\usepackage[margin=2.2cm, top=2.5cm, bottom=2.5cm, headheight=18pt]{geometry}

\usepackage{mathptmx}
\usepackage{setspace}
\setstretch{1.15}

\usepackage{graphicx}
\graphicspath{{./}{imagenes/}{./imagenes/}}
\usepackage{float}
\usepackage{booktabs}
\usepackage{tabularx}
\usepackage{array}
\usepackage{multirow}
\usepackage{caption}
\captionsetup{font=small, labelfont=bf, skip=6pt}

\usepackage[table]{xcolor}

\usepackage{enumitem}
\setlist{topsep=2pt, itemsep=2pt, parsep=0pt}

\usepackage{fancyhdr}
\pagestyle{fancy}
\fancyhf{}
\fancyhead[L]{\small\itshape UNSAAC -- Ingeniería de Software I}
\fancyhead[R]{\small\itshape Lab 05: Planning Poker -- SIGC-CUSCO}
\fancyfoot[C]{\thepage}
\renewcommand{\headrulewidth}{0.4pt}
\renewcommand{\footrulewidth}{0pt}

\usepackage{xurl}
\usepackage{hyperref}
\hypersetup{
    colorlinks=true,
    breaklinks=true,
    linkcolor=blue!75!black,
    urlcolor=blue!75!black,
    citecolor=blue!75!black,
    pdftitle={Lab 05 - Estimación de Esfuerzo con Planning Poker - SIGC-CUSCO},
    pdfauthor={Choquenaira, Yaranga, Ccama}
}

\usepackage{microtype}
\emergencystretch=2.5em

\setlength{\parindent}{0.6cm}
\setlength{\parskip}{0.28em}

% Tipos de columnas convenientes para tablas de requerimientos y backlog
\newcolumntype{Y}{>{\raggedright\arraybackslash}X}
\newcolumntype{P}[1]{>{\raggedright\arraybackslash}p{#1}}
\newcolumntype{C}[1]{>{\centering\arraybackslash}p{#1}}

% Colores institucionales y de clasificación ágil
\definecolor{headerblue}{RGB}{20, 65, 110}
\definecolor{rowalt}{RGB}{248, 250, 252}
\definecolor{cortaGreen}{RGB}{220, 245, 230}
\definecolor{cortaText}{RGB}{22, 101, 52}
\definecolor{mediaOrange}{RGB}{254, 243, 199}
\definecolor{mediaText}{RGB}{180, 83, 9}
\definecolor{complejaRed}{RGB}{254, 226, 226}
\definecolor{complejaText}{RGB}{185, 28, 28}

% Comandos auxiliares para etiquetas visuales de tamaño
\newcommand{\badgeCorta}{\colorbox{cortaGreen}{\textcolor{cortaText}{\textbf{\scriptsize CORTA (1--3)}}}}
\newcommand{\badgeMedia}{\colorbox{mediaOrange}{\textcolor{mediaText}{\textbf{\scriptsize MEDIA (5--8)}}}}
\newcommand{\badgeCompleja}{\colorbox{complejaRed}{\textcolor{complejaText}{\textbf{\scriptsize COMPLEJA (13+)}}}}
